\documentclass[a4paper,11pt]{article}

\usepackage[utf8]{inputenc}

\usepackage{minted}

\usepackage{graphicx}
\usepackage{caption}
\usepackage{subcaption}
\usepackage{placeins}
\usepackage{pgfplots}
\pgfplotsset{compat=1.18} 

\begin{document}

\title{
    \textbf{Reduce Part 2}
}
\author{Jonathan Karlsson}
\date{Spring Term 2024}

\maketitle

\section*{Introduction}
This task is a continuation of the previous part where the three functions map, reduce and filter should be implemented.
The functions from the previous part should also be rewritten using these functions. A new function for taking a list of integers and returning a sum of the square of all values less than a specified value should also be created. 

\section*{Implementation}

The first function that was created was the map function, this function simply copies all elements of the inputted list to a new list, but first runs each element through a function given as an argument.
\begin{minted}{elixir}
def map(list, func) do
  map(list,func,[])
end

def map([h|t],func,new_list) do
  map(t,func,new_list++[func.(h)])
end

def map([],_,new_list) do
  new_list
end
\end{minted}
The next function was the reduce. It was implemented in two different ways. reducel and reducer. 



\end{document}
